% !PRINTED 484
\weberpage{484}

For if under $\pi$ the coefficient $A_{h,i,k}$ passes into $A_{h,i,k}'$, then by means of the transformation $(y) = A(x)$:
\[
\sum A_{h,i,k}' \, y_h y_i y_k = \sum A_{h,i,k} \, x_h x_i x_k,
\]
and by means of the substitution $\xi = A_1(\eta)$ the identity $f(\xi_1, \xi_2, \xi_3) = f(\eta_1, \eta_2, \eta_3)$ holds. Consequently, by the theorem of \S.\,129, 3., there results:
\[
\sum A_{h,i,k}' \, \frac{\partial^3 f}{\partial\eta_h\partial\eta_i\partial\eta_k} = \sum A_{h,i,k} \, \frac{\partial^3 f}{\partial\xi_h\partial\xi_i\partial\xi_k}.
\]

If, then, $\pi$ denotes any permutation of the group $P_{168}$ by which $X_1, X_2, X_3$ pass into $Y_1, Y_2, Y_3$, and
\[
A = \begin{pmatrix} a_1, & b_1, & c_1\\ a_2, & b_2, & c_2\\ a_3, & b_3, & c_3 \end{pmatrix}
\]
is the corresponding substitution from the group $G_{168}$, then we have to set:
\[
\begin{aligned}
\xi_1 &= a_1\eta_1 + a_2\eta_2 + a_3\eta_3\\
\xi_2 &= b_1\eta_1 + b_2\eta_2 + b_3\eta_3\\
\xi_3 &= c_1\eta_1 + c_2\eta_2 + c_3\eta_3,
\end{aligned}
\]
and the invariant property of $L$ gives the relation
\weq{10}{Y_1\eta_1 + Y_2\eta_2 + Y_3\eta_3 = X_1\xi_1 + X_2\xi_2 + X_3\xi_3,}
or, developed,
\weq{11}{\begin{aligned} Y_1 &= a_1X_1 + b_1X_2 + c_1X_3\\ Y_2 &= a_2X_1 + b_2X_2 + c_2X_3\\ Y_3 &= a_3X_1 + b_3X_2 + c_3X_3, \end{aligned}}
i.\,e. the permutation $\pi$, applied to the functions $X_1, X_2, X_3$, has the same effect as the linear substitution $A(X_1, X_2, X_3)$, and accordingly the $X_1, X_2, X_3$ are functions of the kind our problem demands (end of \S.\,128).

\begin{center}
\S. 131.

\gesperrt{Possibility of the determination of the functions $X_1, X_2, X_3$.}
\end{center}

In order to secure completely the reduction of the equation of degree 7 with the group $P_{168}$ to the form problem of the group $G_{168}$, one thing still remains:
%CONT>>
